\subsection{Twist-resolved subcategory constraints at higher multiplicity}
\label{subsec:round4-higher}

The projection in Proposition~\ref{prop:central-subcategory-projection}
can be used before enumerating a complete modular datum.  In particular,
it couples a restriction column to its twist; checking the column and the
twist separately loses information.  We record two explicit source tables
which make this necessary condition fully effective.

\paragraph{The half-$E_6$ subcategory.}
Suppose that the based subring on $\{1,g,x\}$ satisfies
$g^2=1$, $gx=xg=x$, and $x^2=1+g+2x$.
Put $d=1+\sqrt3$ and $D_0=2+d^2=6+2\sqrt3$.
Hagge--Hong classify its four categorifications
\cite[Theorem~1 and Sections~4.7--4.8]{HH3}.
The relations imply that every tensor automorphism of the identity is
trivial: its scalar on $x$ satisfies $\eta_x^2=\eta_x$, and then
$\eta_g=1$.  Hence a pivotal structure, when present, is unique.
Two Galois conjugates have $\dim(x)=1-\sqrt3<0$ and cannot be
pseudo-unitary.  The two remaining cases are complex conjugate.
Their center restriction columns and twists are given by
Hong--Rowell--Wang \cite[Section~2.1]{HRWexotic}.  Writing $\zeta=\exp(2\pi i/12)$, one table is
\begingroup\setlength{\arraycolsep}{3pt}
\[
\begin{array}{c|cccccccccc}
Z&1&Y&X_1&X_2&X_3&X_4&X_5&U&V&W\\ \hline
F(Z)&1&g&x&x&x&x&x&1+x&g+x&1+g+x\\
\theta_Z&1&-1&\zeta^{-3}&\zeta^5&\zeta^5&\zeta^2&\zeta^{-4}&1&-1&1.
\end{array}
\]
\endgroup
The other table conjugates every twist.  Complex conjugation of the entire
putative pseudo-unitary category preserves its based ring and its positive
dimensions, and exchanges these two subcategory choices.  Thus exclusion
for the displayed choice excludes both choices; this is a reduction of
existence, not an omission of a categorical branch.

For a high induction column $c$, the unit coordinate is zero.
Consequently, $c_g>0$ forces $\theta=-1$ and $c_g\ge c_x$;
if $c_g=0<c_x$, the twist must be one of the four displayed $X$-twists.
A low column has twist one, and its coordinates must satisfy
$0\le c_g\le c_x\le c_1$.
These statements concern all structure constants, with the subring
located by its multiplication table, rather than by census labels.

The twist masses give additional equalities.  If the ambient category has
FP dimension $D$, Corollary~\ref{cor:subcategory-twist-masses} implies
\[
 \sum_{Z:\,\theta_Z=\lambda} c_{x,Z}\operatorname{FPdim}(Z)
       =\frac{D}{D_0}M_\lambda,
\]
where the successive values of $M_\lambda$ at
$\lambda=\zeta^{-3},\zeta^5,\zeta^2,\zeta^{-4},-1,1$
are
\[
 d,\quad 2d,\quad d,\quad d,\quad d+1,\quad 2d+3.
\]
Known low columns are subtracted from the right-hand sides.
The existence reduction by conjugating the entire category, given above,
covers the conjugate source table as well.  The implementation combines it with the full
induction Gram equations, the integrality of joint character spectra,
and equality of multiplicities along Galois orbits.  Optional refinements
add the low Verlinde coefficients and restriction products.

\paragraph{Cyclic pointed subcategories of order three.}
For a generator $g$ of a pointed $C_3$ subcategory, all associator classes
are represented by $s\in\{0,1,2\}$.  A standard cocycle representative is
\[
 \omega_s(a,b,c)=\exp\!\left(\frac{2\pi i s a}{3}
                         \left\lfloor\frac{b+c}{3}\right\rfloor\right),
 \qquad 0\le a,b,c<3.
\]
Solving the one-dimensional half-braiding equations gives, over the
underlying object $g^a$, the three twists
\[
 \exp\!\left(2\pi i\left(\frac{sa^2}{9}+\frac{ab}{3}\right)\right),
 \qquad b\in\{0,1,2\}.
\]
Thus a central object whose restriction contains $g$ or $g^2$ has
$\theta^3=\exp(2\pi i s/3)$.  For the ambient field
$\mathbb Q(\zeta_{24})$, the three height-three relative-trace types
when $s=0$ are $(\lambda,\lambda^2,1)$, $\lambda^3=1$.
For $s=1,2$, the three possible ninth roots have the same normalized
relative-trace type $(0,0,\exp(2\pi i s/3))$.
All three associator classes must be retained.  No triviality assumption
on the pointed associator is used.

\paragraph{Exact arithmetic and scope.}
The supplementary implementations retain rational equations and integer
separating rays.  For the half-$E_6$ input with FP field
$\mathbb Q(\sqrt3,\sqrt5)$, normalized traces from
$\mathbb Q(\zeta_{60})$ to the real FP field reduce the height-three
twist domain from $131$ to $40$ types.  A field pullback is accepted only
after exact reconstruction of every coordinate.  The Galois group of
this real field has exponent two, so the square action on these relative
twist traces is trivial; joint-spectrum multiplicity equalities can
therefore be imposed separately for each relative-trace type.
Additional negative-stabilizer rotation conditions are imposed only
after exhaustive reconstruction of the relevant involution module.
Surviving systems remain necessary relaxations and do not provide
categorifications or modular data.


\paragraph{Certified applications and the remaining multiplicity-five input.}
For $R_{6,6,49}$ the complete low domain has $353$ frames, of which only
one remains after the ordinary and character-assignment checks.  Its four assignments, three pointed
associator classes, and two second-indicator choices give $24$ exact
separating rays.  The established two-FP-value reduction gives unrestricted
categorical scope.
For $R_{6,6,1035}$ there are $12$ low frames, of which two admit a character
assignment.  Each has two second-indicator choices; all four full
cyclotomic projection systems have exact separating rays.  The unrestricted
signed reduction exhausts $448$ signed tables in $14$ sign orbits and
reduces every possible categorification to the pseudo-unitary problem.
These two exclusions complete the exact multiplicity-six stratum.
The final half-$E_6$ certificate uses the full conductor-$60$ system;
its validity does not depend on the exploratory real-field MILP
infeasibility results or on additional Clifford rotation constraints.

The multiplicity-five input $R_{6,5,1204}$ remains unresolved.  Its
$39$ low frames admit eight character assignments across two frames.
The height-three full-character system, and the subsequent low-Verlinde
and restriction-product refinement, exclude the two assignments of the
first frame but retain a necessary relaxation on the second.  Neither
numerical feasibility nor the failure of particular Morita model searches
is a categorification or a general exclusion.
